%------------------------------------------------------------------------------%
\section{Revisão}

\begin{exercises*}

%------------------------------------------------------------------------------%
\begin{exercise}
Calcule as Integrais, usando as técnicas estudadas em sala.
\begin{mathlist}
  \item \int x \cos(5x)dx
  \item \int p^5 \ln(p) dp
  \item \int \arccos(x) dx
  \item \int_4^9 \frac{\ln(y)}{\sqrt{y}} dy
  \item \int_0^1 (x^2+1)e^{-x}dx
  \item \int x \sin(x^2) dx
  \item \int \frac{(\ln x)^2}{x} dx
  \item \int e^{\tan(x)} \sec^2(x)dx
  \item \int \frac{dx}{\sqrt{1-x^2} \arcsin(x)}
  \item \int \sin(t) \sec^2(\cos(t))dt
  \item \int x^3\sqrt{x^2+1} dx
  \item \int_e^{e^4} \frac{dx}{x \sqrt{\ln(x)}}
  \item \int_0^1 (3t-1)^{50} dt
  \item \int \sin^6(x)\cos^3(x)dx
  \item \int_0^{\pi} cos^4(2t) dt
  \item \int  \cos(\theta) \cos^5(\sin(\theta)) d \theta
  \item \int_0^1 x^3 \sqrt{1-x^2} dx
  \item \int_{\sqrt{2}}^2 \frac{1}{t^3(\sqrt{t^2-1})} dt
  \item \int \frac{dx}{\sqrt{x^2+16}} dx
  \item \int_0^{a} x^2 \sqrt{a^2-x^2}dx
  \item \int \frac{x-9}{(x+5)(x-2)} dx
  \item \int_0^1 \frac{2}{2x^2+3x+1} dx
  \item \int_3^4 \frac{x^3-2x^2-4}{x^3+2x^2} dx
  \item \int \frac{x^2+1}{(x-3)(x-2)^2} dx
  \item \int \frac{10}{(x-2)(x^2+9)} dx
\end{mathlist}
\begin{answer}
\begin{mathlist}
  \item \frac{1}{5}x \sin(5x)+\frac{1}{25} \cos(5x)+c
  \item \frac{1}{6}p^6 \ln(p) +\frac{1}{36}p^6+C
  \item ~
  \item 6\ln(9)-4\ln(4)-4
  \item -\frac{6}{e}+3
  \item -\frac{1}{2} \cos(x^2)+C
  \item \frac{1}{2}(\ln(x))^2+C
  \item e^{\tan(x)}+C
  \item \ln(\vert \arcsin(x) \vert)+ C
  \item -\tan(\cos(x))+C
  \item \frac{1}{5}(x^2+1)^{\frac{5}{2}}- \frac{1}{3}(x^2+1)^{\frac{3}{2}}+C
  \item 2
  \item \frac{1}{153}(2^{51}+1)
  \item \frac{1}{7} \sin^7(x)-\frac{1}{5} \sin^5(x)+C
  \item \frac{3\pi}{8}
  \item \sin(\sin(\theta))- \frac{2}{3}\sin^5(\sin(\theta))+ \frac{1}{5} \sin^5(\sin(\theta))+C
  \item \frac{2}{5}
  \item \frac{\pi}{4}+ \frac{\sqrt{3}}{8}-\frac{1}{4}
  \item \ln(\sqrt{x^2+16}+x)-\ln(4)+C
  \item \frac{\pi}{16}a^4
  \item 2 \ln \vert x+5 \vert-\ln \vert x-2 \vert+C
  \item 2\ln(\frac{3}{2})
  \item \frac{7}{6}+ \ln(\frac{2}{3})
  \item 10\ln \vert x-3 \vert-9 \ln \vert x-2 \vert+ \frac{5}{x-2}+C
  \item \frac{3}{8} \pi
\end{mathlist}
\end{answer}
\end{exercise}

%------------------------------------------------------------------------------%
\begin{exercise}
Calcule as Integrais. Em cada um dos itens você poderá de usar mais de uma
técnica ou até a mesma técnica mais de uma vez,
\begin{mathlist}
  \item \int_0^{\pi} t \cos^2(t) dt
  \item \int x \arccos(x)dx
  \item \int \frac{cos^3(\sqrt{x})}{\sqrt{x} \hspace{0,1cm} cossec(\sqrt{x})}dx
  \item \int \frac{e^{-x}}{\sqrt{1-e^{-2x}}}dx
 %\item \int \frac{r}{\sqrt{r^2-2r}} dr
  \item \int x^5e^{-x^3} dx
  \item \int \frac{\sec(\theta)\tan(\theta)}{\sec^2(\theta)- \sec(\theta)}d \theta
  \item \int \frac{\ln(x)}{x \sqrt{1+(\ln(x))^2}}dx
  \item \int \frac{e^x}{(e^{2x}+6e^x+10)^{\frac{3}{2}}} dx
  \item \int \frac{3x^2-2}{x^3-2x-8}dx
  \item \int \frac{(\ln(x))^2}{x^3}dx
  \item \int \frac{\ln(x+1)}{x^2}dx
  \item \int x^2 \arcsin(x)dx
  \item \int \frac{\sin^3(\frac{1}{x})}{x^2}dx
  \item \int \frac{\sin^3(x)\cos(x)}{\sqrt{16+\sin^2(x)}}dx
  \item \int x \arctan(x) dx
  \item \int \frac{\sec^2(t)}{\tan^2(t)+3\tan(t)+2}dt
\end{mathlist}
\end{exercise}

%------------------------------------------------------------------------------%
\begin{exercise}
Uma partícula se move ao longo de uma reta com uma velocidade
$v(t)= t^2-t$, onde $v$ é medida em metros por segundo. Ache
\begin{alphalist}
  \item O deslocamento.
  \item A distância percorrida pela partícula no intervalo de tempo $[0,5]$.
\end{alphalist}
\textit{Lembrete}: Se $s(t)$ é a função deslocamento, então a velocidade
$v(t)=s'(t)$. da mesma forma se $v(t)$ é a equação da velocidade então
a função posição (deslocamento) é $s(t)=\int v(t) dt$, caso venha pedido
a distância percorrida no intervalo $[a,b]$ então \(s= \int_a^b v(t) dt\).
O deslocamento nos primeiros $t$ segundos é dado por \(s(t)= \int_0^t v(t) dt\).
\begin{answer}
\begin{mathlist}
  \item s(t)= \frac{t^3}{3}-\frac{t^2}{2}+C
  \item \frac{175}{6}m/s
\end{mathlist}
\end{answer}
\end{exercise}

%------------------------------------------------------------------------------%
\begin{exercise}
Se $f$ for uma função contínua tal que
\[
  \int_1^x f(t) dt= (x-1)e^{2x}+ \int_1^x e^{-t}f(t) dt
\]
para todo $x$, ache uma formula explícita para $f(x)$.
\hint Use o Teorema Fundamental do cálculo parte 1.
\begin{answer}
  \(f(x)= \frac{e^{2x}(2x-1)}{1-e^{-x}}\)
\end{answer}
\end{exercise}

%------------------------------------------------------------------------------%
\begin{exercise}
  Seja
  \(g(x)=\int_1^x f(t) dx\) e
  \(f(t)=\int_1^{t^2} \frac{\sqrt{1+u^7}}{\sqrt{u^5}}du\),
  encontre $g''(2)$.
\begin{answer}
  \(\frac{\sqrt{1+2^{14}}}{16}\)
\end{answer}
\end{exercise}

%------------------------------------------------------------------------------%
\begin{exercise}
Seja \(f(x)=\int_0^x (1-t^2)e^{t^2} dt\).
\begin{alphalist}
  \item Encontre e classifique os pontos críticos de $f$.
  \item Analise o crescimento e o decrescimento de $f$.
\end{alphalist}
\begin{answer}
\begin{alphalist}
  \item Pontos críticos: $x=-1$ e $x=1$.
  \item $f$ é crescente no intervalo $(-1,1)$ e
        decrescente nos intervalos $(-\infty,-1)$ e $(1, \infty)$.
\end{alphalist}
\end{answer}
\end{exercise}

%------------------------------------------------------------------------------%
\begin{exercise}
Uma partícula se move ao longo de uma reta com velocidade igual a
$v(t)=t^2e^{-t}$ metros por segundo após $t$ segundos.
Qual a distância que essa partícula percorrerá nos primeiros $t$ segundos.
\begin{answer}
  $s(t)=-e^{-t}(t^2+2t+2)+K$
\end{answer}
\end{exercise}

%------------------------------------------------------------------------------%
\begin{exercise}
Uma partícula se move em linha reta com função velocidade
$v(t)= \sin(\omega t) \cos^2(\omega t)$.
Encontre a sua função posição $s=f(t) $ se $f(0)=0$?
\begin{answer}
  \(S=\frac{1-\cos^3(\omega t)}{3\omega}\)
\end{answer}
\end{exercise}

%------------------------------------------------------------------------------%
\begin{exercise}
Use substituição trigonométrica para mostrar que
\[
  \int \frac{dx}{\sqrt{x^2+a^2}}= \ln(x+\sqrt{x^2+a^2})+C
\]
\end{exercise}


%------------------------------------------------------------------------------%
\begin{exercise}
Encontre a área delimitada pela elipse
\[
  \frac{x^2}{a^2}+\frac{y^2}{b^2}=1
\]
\begin{answer}
  $\pi ab$
\end{answer}
\end{exercise}

%------------------------------------------------------------------------------%
\begin{exercise}
Encontre a área sob a curva $y=e^{\sin(x)}\sin(2x)$ de $x=0$ a $x=\frac{\pi}{2}$.

\textit{Lembrete}:
Lembre que a área sob (abaixo) de uma curva $y=f(x)$ (limitada pelo eixo x),
no intervalo $[a,b]$  é dada por \(A=\int_a^b f(x) dx\).
Além disso, lembre-se que $\sin(2x)=2\sin(x)\cos(x)$.
\begin{answer}
  $A=2$
\end{answer}
\end{exercise}

%------------------------------------------------------------------------------%
\begin{exercise}
Se $f$ é uma função contínua no intervalo $[a,b]$,
o valor médio de $f$ é definido como
\[
  f_{\text{med}} = \frac{1}{b-a} \int_a^b f(x) dx
\]
Encontre o valor médio das seguintes funções
\begin{enumerate}
  \item $f(x)=x \sec^2(x)$ no intervalo $[0, \frac{\pi}{4}]$.
  \item $f(x)=\frac{\sqrt{x^2-1}}{x}$, $1 \leq x \leq 7$.
\end{enumerate}
\begin{answer}
\begin{mathlist}
  \item V_m=1-\frac{\ln(4)}{\pi}
  \item \frac{1}{6}(4\sqrt{3}-arcsec(7))
\end{mathlist}
\end{answer}
\end{exercise}

%------------------------------------------------------------------------------%
\begin{exercise}
Nas figuras a  seguir, a área 1 ($A_1$) representa a área sob a curva
$f(x)=x^2 \arcsin(x)$ e a área 2  ($A_2$) representa a área sob a curva
$g(x)=x \arccos(x^2)$. Encontre a diferença $A_2-A_1$.
\textit{OBS}: Para calcular a integral envolvida no cálculo da área
$A_2$, você deve OBRIGATORIAMENTE começar usando uma substituição simples.
\begin{center}
  \includegraphics[width=0.7\linewidth]{exercicio_area_1_1}
  \includegraphics[width=0.7\linewidth]{exercicio_area_1_2}
\end{center}
\end{exercise}

%------------------------------------------------------------------------------%
\begin{exercise}
Encontre a área delimitada pelas curvas $y=1-x^2$ e $y=3x^2$.
\begin{answer}
  \(A= \frac{11}{24}\)
\end{answer}
\end{exercise}

%------------------------------------------------------------------------------%
\begin{exercise}
Encontre a área delimitada pelas curvas $y=3x$ e $y=2x^2+x$.
\begin{answer}
  $A= \sfrac{1}{3}$
\end{answer}
\end{exercise}

%------------------------------------------------------------------------------%
\begin{exercise}
Encontre a área delimitada pelas curvas $y=\sqrt[3]{x}$ e $y=x$.
\begin{answer}
  $A= \sfrac{1}{2}$
\end{answer}
\end{exercise}

%------------------------------------------------------------------------------%
\begin{exercise}
Encontre a área delimitada pelas curvas $y= \cos(x)$ e $y= \cos^2(x)$
entre $x=0$ e $x=\pi$.
\begin{answer}
  $A= 2$
\end{answer}
\end{exercise}

%------------------------------------------------------------------------------%
\begin{exercise}
Um arquiteto propôs colocar no piso de uma praça pastilhas coloridas em uma
região plana cujo o formato se assemelha (respeitando a escala) a região
delimitada simutâneamente pelas quarto curvas $y= \ln(x)$, $y=e^x$,
$y=1-x$ e $x=e$. Esboce a região e calcule a área.
\end{exercise}

%------------------------------------------------------------------------------%
\begin{exercise}
Determine a área da região limitada exatamente pelas três curvas
$y =x^2$, $ y = 2- x^2$ e $y = 2x + 8$.
\begin{answer}
  $\frac{100}{3}$
\end{answer}
\end{exercise}

%------------------------------------------------------------------------------%
\begin{exercise}
Avalie as seguintes integrais impróprias:
\begin{mathlist}
  \item \int_{1}^{\infty} \frac{1}{(3x+1)^2}dx
  \item \int_{-\infty}^{\infty} x^3e^{-x^4} dx
  \item \int_{-1}^0 \frac{e^{\frac{1}{x}}}{x^3}dx
  \item \int_3^5 \frac{dx}{x^2-5x+6}
  \item \int_{-2}^{3} \frac{1}{x^4} dx
\end{mathlist}
\begin{answer}
\begin{alphalist}
  \item $\sfrac{1}{12}$
  \item 0
  \item $\sfrac{-2}{e}$.
  \item A integral é divergente.
  \item A integral é divergente.
\end{alphalist}
\end{answer}
\end{exercise}

%------------------------------------------------------------------------------%
\begin{exercise}
A área sob a curva $y=xe^{-x}$ definida no intervalo $[0,\infty)$ é convergente?
\begin{answer}
  A área converge e vale 1
\end{answer}
\end{exercise}

%------------------------------------------------------------------------------%
\begin{exercise}
Podemos estender a nossa definição de valor médio de uma função contínua
a um intervalo definindo o valor médio de $f$ no intervalo $[a, \infty)$
como
\[
  \lim_{ t \to \infty} \frac{1}{t-a} \int_a^t f(x) dx
\]
Calcule o valor médio de $y=\arctan(x)$ no intervalo $[0, \infty)$.
\begin{answer}
  $A= \frac{\pi}{2}$
\end{answer}
\end{exercise}

%------------------------------------------------------------------------------%
%\begin{exercise}
%Se $f$ for uma função contínua tal que
%$$ \int_1^{x^2} f(t) dt= x\arctan(e^x)+ \int_1^x e^{x}f(t) dt$$
%para todo $x$, ache uma formula explícita para $f(x) $.
%\begin{answer}
%\end{answer}
%\end{exercise}

%------------------------------------------------------------------------------%
\begin{exercise}
Avalie a integral imprópria \(\int_0^{\infty} \frac{e^{-x}}{1-e^{-2x}}dx\).
Lembre-se que a integral indefinida já foi resolvida no exercício 2-d.
\begin{answer}
  A integral é divergente.
\end{answer}
\end{exercise}

%------------------------------------------------------------------------------%
%\begin{exercise}
% Seja $g(x)=\int_1^x f(t) dx$ e $f(t)=\int_1^{t^2} \frac{\sqrt{1+u^7}}{\sqrt{u^5}}du$, encontre $g''(2)$.
%\item Determine os pontos de inflexão da curva $f(x)= \int_0^x \frac{1}{3+2t+4t^2}$.
%\item Calcule $\displaystyle \lim_{x \to 3} \left( \frac{x}{x-3} \int_3^x \frac{\sin(t)}{t}dt\right)$.
%\begin{answer}
%\end{answer}
%\end{exercise}

%------------------------------------------------------------------------------%
%\begin{exercise}
%\begin{answer}
%\end{answer}
%\end{exercise}

%------------------------------------------------------------------------------%
%\begin{exercise}
%\begin{mathlist}[2]
%  \item
%  \item
%  \item
%\end{mathlist}
%\begin{answer}
%  \begin{alphalist}[3]
%    \item
%    \item
%    \item
%  \end{alphalist}
%\end{answer}
%\end{exercise}

\end{exercises*}
%------------------------------------------------------------------------------%
